% Preamble template for Tufte-style applied econometrics lecture notes.
% Do not compile this file directly; compile the main notes file instead.
\documentclass{tufte-handout}

\usepackage[T1]{fontenc}
\usepackage[utf8]{inputenc}
\usepackage{ntheorem}
\usepackage{graphicx}
\usepackage{amsmath}
\usepackage{amssymb}
\usepackage{mathtools}
\usepackage{hyperref}
\usepackage{epigraph}
\usepackage{booktabs}
\usepackage{array}
\theoremstyle{break}

\hypersetup{
  colorlinks,
  urlcolor = blue,
  pdfauthor={Swapnil Singh},
  pdfkeywords={econometrics, causal inference, applied econometrics},
  pdftitle={Lecture Notes for Applied Econometrics},
  pdfpagemode=UseNone
}

\newtheorem{ruleN}{Rule}
\newtheorem{thmN}{Theorem}
\newtheorem{assN}{Assumption}
\newtheorem{defN}{Definition}
\newtheorem{exmp}{Example}
\newtheorem{cmt}{Comment}
\newtheorem{proof}{Proof}

\newcommand{\continuation}{??}
\newtheorem*{excont}{Example \continuation}
\newenvironment{continueexample}[1]
 {\renewcommand{\continuation}{\ref{#1}}\excont[continued]}
 {\endexcont}

\newcommand{\bY}{\mathbf{Y}}
\newcommand{\bX}{\mathbf{X}}
\newcommand{\bD}{\mathbf{D}}
\newcommand{\E}{\mathbb{E}}
\newcommand{\Prob}{\mathbb{P}}
\newcommand{\Var}{\mathrm{Var}}
\newcommand{\Cov}{\mathrm{Cov}}
\newcommand{\se}{\mathrm{SE}}
\newcommand{\CI}{\mathrm{CI}}
\newcommand{\ATE}{\mathrm{ATE}}
\newcommand{\ATT}{\mathrm{ATT}}
\newcommand{\TOT}{\mathrm{TOT}}
\newcommand{\LATE}{\mathrm{LATE}}
\newcommand{\ITT}{\mathrm{ITT}}
\newcommand{\RD}{\mathrm{RD}}
\newcommand{\DiD}{\mathrm{DiD}}
\newcommand{\MMSE}{\mathrm{MMSE}}
\newcommand{\OLS}{\mathrm{OLS}}
\newcommand{\Normal}{\mathcal{N}}
\newcommand{\plim}{\operatorname*{plim}}
\DeclareMathOperator*{\argmin}{arg\,min}
\DeclareMathOperator{\Supp}{Supp}

\newcommand\independent{\protect\mathpalette{\protect\independenT}{\perp}}
\def\independenT#1#2{\mathrel{\rlap{$#1#2$}\mkern2mu{#1#2}}}
\newcommand{\indep}{\perp\!\!\!\perp}

\usepackage{cleveref}
\crefname{appsec}{appendix}{appendices}
\crefname{appsubsec}{appendix}{appendices}
\crefname{assumption}{assumption}{assumptions}
\crefname{equation}{equation}{equations}
\crefname{exmp}{example}{examples}
\crefname{assN}{assumption}{assumptions}
\crefname{cmt}{comment}{comments}
\crefname{defN}{definition}{definitions}

\usepackage[nolist]{acronym}
\begin{acronym}
  \acro{ATE}{average treatment effect}
  \acro{ATT}{average treatment effect on the treated}
  \acro{CIA}{conditional independence assumption}
  \acro{CI}{confidence interval}
  \acro{CLT}{central limit theorem}
  \acro{DiD}{difference-in-differences}
  \acro{FE}{fixed effects}
  \acro{FWL}{Frisch-Waugh-Lovell}
  \acro{IV}{instrumental variables}
  \acro{LATE}{local average treatment effect}
  \acro{OLS}{ordinary least squares}
  \acro{OVB}{omitted-variable bias}
  \acro{PO}{potential outcomes}
  \acro{RCT}{randomized controlled trial}
  \acro{RD}{regression discontinuity}
  \acro{SUTVA}{stable unit treatment value assignment}
  \acro{TWFE}{two-way fixed effects}
\end{acronym}

\def\inprobHIGH{\,{\buildrel p \over \rightarrow}\,}
\def\inprob{\,{\inprobHIGH}\,}
\def\indistHIGH{\,{\buildrel d \over \rightarrow}\,}
\def\indist{\,{\indistHIGH}\,}

\usepackage[many]{tcolorbox}
\definecolor{main}{HTML}{3F6EA7}
\definecolor{sub}{HTML}{F2F7FD}
\definecolor{sub2}{HTML}{FFF5E8}
\definecolor{mainline}{HTML}{D3DFEE}
\definecolor{warn}{HTML}{C7771B}
\definecolor{warnline}{HTML}{EAD9C0}

\tcbset{
    enhanced,
    breakable,
    colback = white,
    boxrule = 0.5pt,
    arc = 2pt,
    left = 9pt,
    right = 9pt,
    top = 7pt,
    bottom = 7pt,
    before skip = 0.3cm,
    after skip = 0.45cm
}

\newtcolorbox{boxD}{
    colback = sub,
    colframe = mainline,
    borderline west = {2.2pt}{0pt}{main},
    borderline north = {0.6pt}{0pt}{main!25},
    borderline south = {0.6pt}{0pt}{main!15}
}

\newtcolorbox{boxF}{
    colback = sub2,
    colframe = warnline,
    borderline west = {2.2pt}{0pt}{warn},
    borderline north = {0.6pt}{0pt}{warn!25},
    borderline south = {0.6pt}{0pt}{warn!15}
}

\usepackage{colortbl}
\usepackage{tikz}
\usepackage{verbatim}
\usetikzlibrary{positioning}
\usetikzlibrary{snakes}
\usetikzlibrary{calc}
\usetikzlibrary{arrows}
\usetikzlibrary{decorations.markings}
\usetikzlibrary{shapes.misc}
\usetikzlibrary{matrix,shapes,arrows,fit,tikzmark}


\title{Applied Econometrics: Session 7 Notes}
\author{Swapnil Singh}
\date{}

\hypersetup{
  pdftitle={Applied Econometrics: Session 7 Notes},
  pdfauthor={Swapnil Singh},
  pdfkeywords={difference-in-differences, panel data, fixed effects, event study, parallel trends}
}

\setlength{\headheight}{26pt}

\begin{document}
\maketitle

\begin{abstract}
This session studies designs that use variation across units and across time.
Difference-in-differences compares changes in treated units with changes in
untreated units. Panel fixed effects compare a unit to itself over time and
remove time-invariant unobserved heterogeneity. The central identifying
assumption is not that treated and control units have the same level of the
outcome. It is that, absent treatment, their outcome changes would have followed
parallel paths. We develop the 2-by-2 estimator, its regression version, event
study diagnostics, one-way and two-way fixed effects, standard errors, and
practical reporting.
\end{abstract}

\begin{boxD}
\textbf{Reading and objective.} The main readings are Chapter 13,
``Difference-in-Differences,'' and Chapter 14, ``Panel Data and Fixed Effects,''
from Matheus Facure Alves, \emph{Causal Inference for the Brave and True}. By
the end of the session, students should be able to compute a 2-by-2 DiD by hand,
state the parallel-trends assumption in potential-outcomes notation, estimate
the same estimand by regression, explain fixed effects as demeaning, interpret a
simple event study, choose reasonable standard errors, and report the design in
a way that a skeptical reader can audit.
\end{boxD}

\section{Why This Lecture Matters}

Many empirical questions are about policy changes, market reforms, school
programs, firm decisions, or local shocks. We rarely get a clean randomized
experiment. Instead, one group is exposed to a change and another group is not.
We observe both groups before and after the change. Difference-in-differences
(DiD) is the simplest way to turn this structure into a causal estimate.

The important word is \emph{difference}. A single before-after comparison is
usually not credible because many outcomes trend over time. A single treated-
control comparison after treatment is usually not credible because groups can
have different baseline levels. DiD removes both problems only if the remaining
change in the control group is a good proxy for the change the treated group
would have experienced without treatment.

\begin{boxF}
\textbf{Core intuition.} DiD does not ask whether treated and control units are
identical. It asks whether the control group supplies the right missing trend
for the treated group. Fixed effects extend the same logic: compare units to
themselves so that stable differences across units cannot explain the estimate.
\end{boxF}

\section{The 2-by-2 DiD Design}

\subsection{The Four Cells}

The cleanest DiD design has two groups and two periods. Let $D_i=1$ denote a
treated group and $D_i=0$ a control group. Let $Post_t=1$ denote the period after
the policy and $Post_t=0$ the period before the policy. There are four observed
cell means:
\[
\begin{array}{c@{\quad}cc}
& Post_t=0 & Post_t=1 \\
\midrule
D_i=0 & \bar{Y}_{00} & \bar{Y}_{01}\\
D_i=1 & \bar{Y}_{10} & \bar{Y}_{11}
\end{array}
\]
The first subscript denotes group and the second denotes time. Thus
$\bar{Y}_{11}$ is the treated group after treatment and $\bar{Y}_{10}$ is the
treated group before treatment.

\begin{defN}[2-by-2 DiD setting]
A 2-by-2 DiD design has one treated group, one comparison group, one pre period,
and one post period. Treatment occurs only for the treated group in the post
period, so the treatment indicator is
\[
    T_{it}=D_i Post_t.
\]
\end{defN}

\subsection{Two Naive Comparisons}

Before writing the DiD estimator, it is useful to see why the two most obvious
estimators are weak.

\paragraph{Naive before-after comparison.}
The treated group's change is
\[
    \bar{Y}_{11}-\bar{Y}_{10}.
\]
This compares treated units after the policy to the same group before the
policy. It controls for permanent differences between treated and control units
because it uses only the treated group. But it does not control for time shocks.
If the outcome would have risen anyway, this comparison exaggerates the policy
effect. If the outcome would have fallen anyway, it understates the effect.

\paragraph{Naive post-period cross-section.}
The post-treatment treated-control gap is
\[
    \bar{Y}_{11}-\bar{Y}_{01}.
\]
This uses the post period only. It controls for common post-period shocks because
both groups are measured at the same time. But it does not control for permanent
baseline differences. If treated units always had higher outcomes, the estimator
will attribute that old difference to the new policy.

\begin{boxD}
\textbf{Teaching point.} Students often say ``we control for the before period''
or ``we control for the control group.'' That is too vague. The before-after
comparison controls for fixed group differences but not time shocks. The
post-only comparison controls for time shocks but not fixed group differences.
DiD combines both comparisons.
\end{boxD}

\subsection{The DiD Estimator}

The DiD estimator subtracts the control group's change from the treated group's
change:
\[
    \widehat{\delta}_{DiD}
    =
    (\bar{Y}_{11}-\bar{Y}_{10})
    -
    (\bar{Y}_{01}-\bar{Y}_{00}).
\]
Equivalently, it subtracts the pre-treatment treated-control gap from the
post-treatment treated-control gap:
\[
    \widehat{\delta}_{DiD}
    =
    (\bar{Y}_{11}-\bar{Y}_{01})
    -
    (\bar{Y}_{10}-\bar{Y}_{00}).
\]
Both formulas are the same algebraically. The first emphasizes changes over
time. The second emphasizes changes in group gaps.

\begin{exmp}[Hand calculation]
Suppose an employment program is introduced in treated cities after period 0.
The outcome is average monthly earnings.
\[
\begin{array}{lcc}
\toprule
& \text{Before} & \text{After}\\
\midrule
\text{Control cities} & 100 & 112\\
\text{Treated cities} & 80 & 105\\
\bottomrule
\end{array}
\]
The treated change is $105-80=25$. The control change is $112-100=12$.
Therefore
\[
    \widehat{\delta}_{DiD}=25-12=13.
\]
The estimate says that earnings rose by 13 more units in treated cities than
would be predicted by the control-city trend.
\end{exmp}

Notice the role of the control group. It is not used because the control cities
have the same level as treated cities. They do not: controls start at 100 and
treated cities start at 80. The control group is used because its increase of 12
is assumed to be the right counterfactual increase for treated cities.

\section{Potential Outcomes and Parallel Trends}

\subsection{Potential Outcomes}

Let $Y_{it}(1)$ be the outcome for unit $i$ at time $t$ if treated and
$Y_{it}(0)$ the outcome if untreated. In the post period, treated units reveal
$Y_{i1}(1)$, but the causal effect for treated units requires the missing
counterfactual $Y_{i1}(0)$:
\[
    \ATT
    =
    \E[Y_{i1}(1)-Y_{i1}(0)\mid D_i=1].
\]
The problem is that for treated units after treatment we observe
$Y_{i1}=Y_{i1}(1)$, not $Y_{i1}(0)$.

\begin{defN}[Observed outcome in 2-by-2 DiD]
With treatment $T_{it}=D_i Post_t$, the observed outcome is
\[
    Y_{it}=T_{it}Y_{it}(1)+(1-T_{it})Y_{it}(0).
\]
Before treatment, treated and control units are untreated, so their observed
pre-period outcomes are untreated potential outcomes.
\end{defN}

\subsection{Counterfactual Imputation}

DiD imputes the missing treated post-period untreated outcome by taking the
treated pre-period level and adding the control group's untreated trend:
\[
    \widehat{\E[Y_{i1}(0)\mid D_i=1]}
    =
    \E[Y_{i0}\mid D_i=1]
    +
    \{\E[Y_{i1}\mid D_i=0]-\E[Y_{i0}\mid D_i=0]\}.
\]
This is the central idea from the Python Causality Handbook's billboard example:
use the untreated city's growth as the missing growth for the treated city. The
treated and control locations may differ in levels, but the untreated trend must
be comparable.

\begin{assN}[Parallel trends]
In a 2-by-2 design, the untreated potential outcome would have changed by the
same amount in treated and control groups:
\[
    \E[Y_{i1}(0)-Y_{i0}(0)\mid D_i=1]
    =
    \E[Y_{i1}(0)-Y_{i0}(0)\mid D_i=0].
\]
\end{assN}

\begin{thmN}[Identification of ATT by DiD]
If treatment affects only the treated group in the post period and parallel
trends holds, then
\[
    \E[\widehat{\delta}_{DiD}] = \ATT.
\]
\end{thmN}

\begin{proof}
Write the population DiD:
\[
\delta =
\{\E[Y_{i1}\mid D_i=1]-\E[Y_{i0}\mid D_i=1]\}
-
\{\E[Y_{i1}\mid D_i=0]-\E[Y_{i0}\mid D_i=0]\}.
\]
For treated units in the post period, $Y_{i1}=Y_{i1}(1)$. For treated units in
the pre period and for all control observations, observed outcomes are untreated
potential outcomes. Therefore
\[
\delta =
\{\E[Y_{i1}(1)\mid D_i=1]-\E[Y_{i0}(0)\mid D_i=1]\}
-
\{\E[Y_{i1}(0)\mid D_i=0]-\E[Y_{i0}(0)\mid D_i=0]\}.
\]
Add and subtract $\E[Y_{i1}(0)\mid D_i=1]$:
\[
\begin{aligned}
\delta
&=
\E[Y_{i1}(1)-Y_{i1}(0)\mid D_i=1]\\
&\quad+
\{\E[Y_{i1}(0)-Y_{i0}(0)\mid D_i=1]
-
\E[Y_{i1}(0)-Y_{i0}(0)\mid D_i=0]\}.
\end{aligned}
\]
The first term is $\ATT$. The second term is zero under parallel trends. Hence
$\delta=\ATT$.
\end{proof}

\subsection{What Breaks Parallel Trends?}

Parallel trends fails when treatment assignment is related to the untreated
growth of the outcome. Examples:
\begin{itemize}
    \item A job-training program is targeted to regions already recovering faster
    than other regions.
    \item A price regulation is introduced in markets where prices were already
    rising unusually fast.
    \item A school intervention is assigned to schools whose test scores were
    already declining.
    \item A firm launches advertising in cities where demand forecasts were
    already improving.
\end{itemize}

The failure can bias estimates upward or downward. If treated units would have
grown faster even without treatment, DiD overstates a positive treatment effect.
If treated units would have grown more slowly even without treatment, DiD
understates the effect.

\section{Regression DiD}

\subsection{The Saturated 2-by-2 Regression}

The four-cell DiD can be estimated by OLS:
\[
    Y_{it}
    =
    \alpha
    + \beta D_i
    + \lambda Post_t
    + \delta(D_i Post_t)
    + u_{it}.
\]
The coefficient $\delta$ equals the DiD estimator.

\begin{defN}[Regression DiD coefficient]
In the saturated 2-by-2 regression, the interaction coefficient on
$D_i Post_t$ is the difference-in-differences estimate.
\end{defN}

To see this, compute the fitted value in each cell:
\[
\begin{array}{c@{\quad}cc}
& Post_t=0 & Post_t=1\\
\midrule
D_i=0 & \alpha & \alpha+\lambda\\
D_i=1 & \alpha+\beta & \alpha+\beta+\lambda+\delta
\end{array}
\]
The treated change is
\[
    (\alpha+\beta+\lambda+\delta)-(\alpha+\beta)=\lambda+\delta.
\]
The control change is
\[
    (\alpha+\lambda)-\alpha=\lambda.
\]
The difference is $\delta$.

\begin{boxF}
\textbf{Why regression is convenient.} Regression gives the same point estimate
in the simple 2-by-2 case, but it also gives standard errors, allows covariates,
handles unbalanced cell sizes, and extends naturally to many units and many
periods.
\end{boxF}

\subsection{Covariates}

A common applied specification is
\[
    Y_{it}
    =
    \alpha
    + \beta D_i
    + \lambda Post_t
    + \delta(D_i Post_t)
    + X_{it}'\theta
    + u_{it}.
\]
Covariates can improve precision and adjust for changes in observed
composition. But they do not rescue a design with bad parallel trends unless
parallel trends becomes plausible \emph{conditional on those covariates}. Be
especially careful with variables that are themselves affected by treatment.
Controlling for post-treatment mediators changes the estimand and can introduce
bias.

\section{Event-Study Intuition}

\subsection{Why More Pre-Periods Help}

With only one pre period, parallel trends cannot be examined visually. With
multiple pre periods, we can plot treated and control outcomes before treatment.
This does not prove the assumption, because the assumption is about the
unobserved post-treatment counterfactual. But it can reveal obvious violations.

\begin{defN}[Event time]
Event time indexes observations relative to the treatment date. If treatment
starts at time $G_i$ for unit $i$, event time is
\[
    k=t-G_i.
\]
Then $k=-1$ is the period immediately before treatment, $k=0$ is the first
treatment period, and $k=1$ is one period after treatment begins.
\end{defN}

\subsection{A Simple Event-Study Regression}

For one treated cohort and a never-treated comparison group, a common event-study
specification is
\[
    Y_{it}
    =
    \alpha_i+\lambda_t+
    \sum_{k\neq -1}\delta_k
    \mathbf{1}\{D_i=1,\ t-G_i=k\}
    +u_{it}.
\]
The unit fixed effects $\alpha_i$ absorb permanent differences across units. The
time fixed effects $\lambda_t$ absorb shocks common to all units in a given
period. The omitted category is usually $k=-1$, the last pre-treatment period.

Pre-treatment coefficients such as $\delta_{-3}$ and $\delta_{-2}$ are called
\emph{leads}. They should be close to zero if treated and control groups were
moving similarly before treatment. Post-treatment coefficients such as
$\delta_0$, $\delta_1$, and $\delta_2$ trace the dynamic treatment effect
relative to the omitted pre-treatment period.

\begin{boxD}
\textbf{How to read an event-study plot.} Look first at the pre-treatment
coefficients. Large systematic pre-trends are a warning sign. Then look at the
post-treatment coefficients: do effects appear immediately, grow gradually,
fade, or reverse? Finally, check the confidence intervals and remember that
failure to reject zero pre-trends is not proof of parallel trends.
\end{boxD}

\subsection{Event-Study Pitfalls}

Event studies are useful, but students should not treat them as magic.
\begin{itemize}
    \item \textbf{Low power.} Pre-trend tests may not reject even when violations
    are economically important.
    \item \textbf{Anticipation.} If units change behavior before official
    treatment, leads may capture real anticipatory effects rather than only bad
    trends.
    \item \textbf{Composition changes.} With staggered adoption, the comparison
    group for each event-time coefficient can change.
    \item \textbf{Bad controls.} Already-treated units can become controls for
    later-treated units in simple TWFE event studies, creating difficult
    weighting problems.
    \item \textbf{Multiple testing.} Many lead and lag coefficients invite
    selective interpretation.
\end{itemize}

This lecture focuses on the simple one-cohort case. In modern applied work with
staggered treatment timing and heterogeneous effects, researchers often use
cohort-specific DiD estimators rather than relying blindly on a single TWFE event
study.

\section{Panel Data}

\subsection{Notation}

Panel data observe the same units repeatedly over time. A unit might be a
person, firm, city, school, bank, country, or product. The panel observation is
indexed by $(i,t)$:
\[
    i=1,\ldots,N,
    \qquad
    t=1,\ldots,T.
\]
A typical panel regression is
\[
    Y_{it}=\alpha+\tau T_{it}+X_{it}'\theta+u_{it}.
\]
The problem is that $u_{it}$ may contain unobserved unit characteristics that
are correlated with treatment. For example, cities with better administrators
may both adopt a policy and have better outcomes.

\begin{defN}[Panel data]
A panel dataset follows the same units over multiple periods. A balanced panel
has every unit observed in every period. An unbalanced panel has some missing
unit-period observations.
\end{defN}

\subsection{Unit Fixed Effects}

Suppose the error contains a time-invariant unit component:
\[
    Y_{it}=\tau T_{it}+X_{it}'\theta+\alpha_i+\varepsilon_{it}.
\]
Here $\alpha_i$ captures permanent differences across units: geography,
baseline ability, long-run institutions, fixed product quality, stable household
preferences, and other features that do not change over the sample.

\begin{defN}[Unit fixed effects]
Unit fixed effects are unit-specific intercepts $\alpha_i$. They allow each unit
to have its own baseline level of the outcome.
\end{defN}

The fixed-effects estimate of $\tau$ is identified from within-unit changes in
treatment. Units that never change treatment status do not directly identify
$\tau$, although they can help estimate time effects and residual variation.

\begin{boxF}
\textbf{Interpretation.} With unit fixed effects, a treated city is compared to
itself before treatment, after removing shocks common to all cities if time
fixed effects are also included. The model does not compare the city level to
another city's level.
\end{boxF}

\subsection{The Within Transformation}

Let
\[
    \bar{Y}_i=\frac{1}{T_i}\sum_t Y_{it},
    \qquad
    \bar{T}_i=\frac{1}{T_i}\sum_t T_{it},
    \qquad
    \bar{X}_i=\frac{1}{T_i}\sum_t X_{it}.
\]
Average the fixed-effects model over time for unit $i$:
\[
    \bar{Y}_i
    =
    \tau \bar{T}_i+\bar{X}_i'\theta+\alpha_i+\bar{\varepsilon}_i.
\]
Subtract this unit average from the original equation:
\[
\begin{aligned}
    Y_{it}-\bar{Y}_i
    &=
    \tau(T_{it}-\bar{T}_i)
    +(X_{it}-\bar{X}_i)'\theta
    +(\alpha_i-\alpha_i)
    +(\varepsilon_{it}-\bar{\varepsilon}_i)\\
    \ddot{Y}_{it}
    &=
    \tau\ddot{T}_{it}
    +\ddot{X}_{it}'\theta
    +\ddot{\varepsilon}_{it}.
\end{aligned}
\]
The unit fixed effect disappears. This is why fixed effects remove
time-invariant confounders.

\begin{thmN}[Within estimator]
OLS on unit-demeaned variables gives the same slope coefficients as OLS with a
full set of unit dummy variables, provided the same observations and regressors
are used.
\end{thmN}

\begin{boxD}
\textbf{Consequence.} Variables that do not change within a unit cannot be
estimated in a unit fixed-effects model. Gender, birth year, permanent location,
or baseline treatment-group status are absorbed by unit fixed effects if they
are constant over time.
\end{boxD}

\subsection{Time Fixed Effects}

Many panels also include shocks that are common to all units in the same period:
inflation, national recessions, pandemics, exchange-rate shocks, common
regulatory changes, or seasonality. Add time fixed effects:
\[
    Y_{it}
    =
    \tau T_{it}+X_{it}'\theta+\alpha_i+\lambda_t+\varepsilon_{it}.
\]
The $\lambda_t$ terms allow each period to have its own intercept. They absorb
common shocks and common trends. Identification now comes from deviations of a
unit from its own average and deviations from the common time pattern.

\begin{defN}[Two-way fixed effects]
A two-way fixed-effects model includes both unit fixed effects $\alpha_i$ and
time fixed effects $\lambda_t$:
\[
    Y_{it}=\tau T_{it}+X_{it}'\theta+\alpha_i+\lambda_t+\varepsilon_{it}.
\]
\end{defN}

\section{TWFE and Simple DiD}

\subsection{Equivalence in the 2-by-2 Case}

In the 2-by-2 design, the regression
\[
    Y_{it}=\alpha+\beta D_i+\lambda Post_t+\delta(D_iPost_t)+u_{it}
\]
is the same as a two-way fixed-effects model with group fixed effects and time
fixed effects:
\[
    Y_{gt}=\alpha_g+\lambda_t+\delta T_{gt}+u_{gt},
\]
where $g$ indexes group and $T_{gt}=D_gPost_t$. The treatment coefficient is the
DiD estimator because the fixed effects remove group levels and period shocks.

\begin{exmp}[TWFE equals DiD in four cells]
Using the earlier table, the DiD estimate was 13. If we run
\[
    Y_{gt}=\alpha_g+\lambda_t+\delta T_{gt}+u_{gt}
\]
on the four cell means, the fitted values imply exactly the same restriction:
the treatment coefficient must explain the extra increase in the treated group
after removing the control-group increase. Hence $\hat{\delta}=13$.
\end{exmp}

\subsection{Two-Way Demeaning}

With unit and time fixed effects, a useful conceptual transformation is
\[
    \tilde{Y}_{it}
    =
    Y_{it}-\bar{Y}_i-\bar{Y}_t+\bar{Y},
\]
and similarly for $T_{it}$ and $X_{it}$. This removes unit means and time means,
while adding back the grand mean because it was subtracted twice. OLS on the
two-way-demeaned variables gives the same slopes as a regression with unit and
time dummies in balanced panels; software implements the corresponding
projection in more general cases.

\section{Standard Errors}

Panel and DiD data usually violate the independent-identically-distributed error
assumption. Outcomes for the same unit are serially correlated. Outcomes within a
state, school, firm, or city may be correlated because the treatment is assigned
at that level.

\begin{boxF}
\textbf{Rule of thumb.} Cluster standard errors at the level of treatment
assignment. If policy varies by state, cluster by state. If treatment varies by
school, cluster by school. If treatment varies by person, cluster by person.
\end{boxF}

Why? Suppose a state policy changes for every individual in the state. Treating
each individual-year observation as independent exaggerates the amount of
information. The number of independent policy shocks is closer to the number of
states than the number of individual observations.

With few clusters, conventional cluster-robust standard errors can still be too
optimistic. Researchers then report wild-cluster bootstrap inference,
randomization inference, or more conservative aggregation. The exact choice
depends on the design, but the principle is constant: inference must respect the
level at which identifying variation occurs.

\section{Diagnostics and Reporting}

\subsection{What to Plot}

Before presenting a DiD estimate, plot the data. At minimum:
\begin{itemize}
    \item group means over time for treated and control groups;
    \item the treatment date marked clearly;
    \item an event-study plot if multiple pre and post periods exist;
    \item sample sizes by group and period;
    \item placebo or falsification outcomes when available.
\end{itemize}

Plots do not identify the model by themselves, but they prevent many bad DiD
applications. If the treated group is already diverging sharply before
treatment, the reader should not be asked to trust a parallel-trends story
without additional evidence.

\subsection{Common Mistakes}

\begin{enumerate}
    \item \textbf{Confusing equal levels with parallel trends.} DiD allows
    different levels. It requires comparable untreated changes.
    \item \textbf{Using only before-after evidence.} A time trend can masquerade
    as a treatment effect.
    \item \textbf{Controlling for post-treatment variables.} Bad controls can
    remove part of the causal effect or open new bias channels.
    \item \textbf{Ignoring treatment timing.} If units adopt treatment at
    different times, simple TWFE may not estimate a transparent average effect.
    \item \textbf{Clustering at the wrong level.} Standard errors should reflect
    the assignment and shock structure.
    \item \textbf{Overclaiming from pre-trend tests.} Not rejecting pre-trends is
    not the same as proving parallel trends.
    \item \textbf{Forgetting the estimand.} Say whether the target is ATT for
    treated units, an average effect over treated unit-periods, or something
    else.
\end{enumerate}

\subsection{Reporting Template}

A clear DiD report should state:
\begin{itemize}
    \item the treated units, control units, treatment date, and sample window;
    \item the outcome definition and measurement frequency;
    \item the exact regression equation;
    \item whether unit fixed effects, time fixed effects, and covariates are
    included;
    \item the identifying assumption in words;
    \item the standard-error clustering level;
    \item the main coefficient, standard error, confidence interval, and units;
    \item graphical evidence on pre-trends and post-treatment dynamics;
    \item robustness checks and limitations.
\end{itemize}

\begin{boxD}
\textbf{Python example.} The file
\texttt{lectures/code/lecture-7-did-panel.py} simulates a panel where parallel
trends holds and a panel where it fails. It computes a 2-by-2 DiD estimate, a
two-way fixed-effects estimate, and event-study coefficients. Use it to see that
the same estimator behaves well when the missing counterfactual trend is valid
and badly when treated units would have grown faster even without treatment.
\end{boxD}

\section{Worked Problems}

\begin{exmp}[Problem 1: Compute DiD from four cells]
A city raises the minimum wage. Average employment is:
\[
\begin{array}{lcc}
\toprule
& \text{Before} & \text{After}\\
\midrule
\text{Control counties} & 52 & 54\\
\text{Treated counties} & 48 & 47\\
\bottomrule
\end{array}
\]
Compute the before-after estimate, the post-only estimate, and the DiD estimate.

\textbf{Solution.} The before-after estimate for treated counties is
$47-48=-1$. The post-only treated-control difference is $47-54=-7$. The DiD
estimate is
\[
    (47-48)-(54-52)=-1-2=-3.
\]
The interpretation is that employment fell by 3 more units in treated counties
than predicted by the control counties' change.
\end{exmp}

\begin{exmp}[Problem 2: State the identifying assumption]
In Problem 1, what must be true for $-3$ to estimate a causal effect?

\textbf{Solution.} Absent the minimum-wage increase, treated counties would have
experienced the same employment change as control counties:
\[
    \E[Y_{1}(0)-Y_{0}(0)\mid D=1]
    =
    \E[Y_{1}(0)-Y_{0}(0)\mid D=0].
\]
In words, the treated counties' employment would have increased by 2 units, just
like the control counties, if the policy had not changed.
\end{exmp}

\begin{exmp}[Problem 3: Fixed-effects demeaning]
Suppose one worker has wages $10,12,14$ over three years and treatment status
$0,1,1$. Compute the within-transformed wage and treatment variables.

\textbf{Solution.} The worker's average wage is $(10+12+14)/3=12$. The average
treatment is $(0+1+1)/3=2/3$. The demeaned wages are $-2,0,2$. The demeaned
treatment values are $-2/3,1/3,1/3$. A unit fixed-effects regression uses this
within-worker variation, not the worker's wage level.
\end{exmp}

\begin{exmp}[Problem 4: Why time fixed effects matter]
All firms in an economy experience an inflation shock in 2025. A subset of firms
also adopts a new technology in 2025. What happens if we omit time fixed effects?

\textbf{Solution.} The technology coefficient may partly capture the economywide
2025 shock. Time fixed effects absorb the common 2025 component, so the treatment
effect is estimated from whether adopting firms changed more than non-adopting
firms in the same year.
\end{exmp}

\section{Checklist for Students}

Before trusting a DiD or fixed-effects estimate, check the following.
\begin{enumerate}
    \item Can I draw the four-cell table or describe the panel structure?
    \item What is the treated group and when does treatment begin?
    \item What is the control group and why is it credible?
    \item What is the exact missing counterfactual?
    \item Is the identifying assumption parallel trends, conditional parallel
    trends, or something else?
    \item Do pre-treatment plots support the story?
    \item Are unit and time fixed effects included for a reason I can explain?
    \item Are covariates pre-treatment or plausibly not affected by treatment?
    \item Are standard errors clustered at the treatment-assignment level?
    \item Does the reported coefficient have a clear unit and estimand?
\end{enumerate}

\section{Practice Questions}

\begin{enumerate}
    \item A tax reform affects one region in 2024. You observe treated and
    control regions in 2022, 2023, 2024, and 2025. Write an event-study
    regression with 2023 as the omitted event-time period.
    \item Explain in words why DiD can allow treated units to have lower baseline
    outcomes than control units.
    \item Give one example of a time-invariant unobserved confounder that unit
    fixed effects can remove and one time-varying confounder that they cannot
    remove.
    \item In the regression $Y_{it}=\alpha_i+\lambda_t+\tau T_{it}+u_{it}$, why
    can a variable such as distance to the capital city not be estimated if it is
    constant within each region?
    \item Suppose treated units have positive and statistically significant
    event-study coefficients two periods before treatment. What does this imply
    for the DiD design?
    \item Why might individual-level observations not justify individual-level
    clustering when treatment is assigned at the state level?
    \item What changes when covariates are added to a DiD regression? What
    assumption is now needed?
    \item Run or read \texttt{lectures/code/lecture-7-did-panel.py}. Compare the
    parallel-trends and nonparallel-trends scenarios. Which coefficients warn
    you that the second design is not credible?
\end{enumerate}

\end{document}
